\documentclass[11pt,a4paper]{article}
\usepackage{fontspec,amsmath,amssymb,amsthm,microtype}
\setmainfont{lmroman10-regular.otf}[BoldFont=lmroman10-bold.otf,ItalicFont=lmroman10-italic.otf,BoldItalicFont=lmroman10-bolditalic.otf]
\usepackage{xcolor}
\usepackage[most]{tcolorbox}
\usepackage{mathrsfs}
\providecommand{\lt}{<}
\providecommand{\gt}{>}
\setmonofont{DejaVuSansMono.ttf}[BoldFont=DejaVuSansMono-Bold.ttf,ItalicFont=DejaVuSansMono-Oblique.ttf,BoldItalicFont=DejaVuSansMono-BoldOblique.ttf,Scale=0.86]
\newfontfamily\leanfallback{STIXTwoMath-Regular.otf}[Scale=0.86]
\definecolor{leanrule}{gray}{0.62}
\definecolor{leanback}{gray}{0.962}
\definecolor{leantitle}{gray}{0.22}
\definecolor{leanmuted}{gray}{0.38}
\newtcolorbox{leanbox}[1]{enhanced,breakable,colback=leanback,colframe=leanrule,
  boxrule=0.4pt,arc=0.6mm,left=2.4mm,right=2.4mm,top=1.8mm,bottom=1.8mm,
  before skip=0.8em,after skip=1em,fonttitle=\small\bfseries,coltitle=leantitle,
  colbacktitle=leanback,titlerule=0.3pt,title={#1}}
\newlength{\leancharwidth}
\newenvironment{leancode}{\par\smallskip\begingroup\ttfamily\footnotesize\raggedright
  \settowidth{\leancharwidth}{M}\setlength{\parskip}{0pt}\setlength{\parindent}{0pt}
  \frenchspacing}{\par\endgroup\smallskip}
\newcommand{\leanline}[2]{\hangindent=\dimexpr#1\leancharwidth+4\leancharwidth\relax
  \hangafter=1\noindent\hspace*{#1\leancharwidth}#2\par}
\newcommand{\leansource}[1]{\par{\footnotesize\color{leanmuted}#1}\par}
\newcommand{\leanlabel}[1]{\par\smallskip\noindent{\small\bfseries #1}\par}
\usepackage[margin=25mm]{geometry}
\usepackage[unicode,hidelinks]{hyperref}
\hypersetup{pdfauthor={LangevinLean Project}}
\setlength{\emergencystretch}{3em}
\setlength{\parskip}{0.4em}
\setlength{\parindent}{0pt}
\allowdisplaybreaks[2]
\raggedbottom
\title{\texorpdfstring{The noise-driven iterates}{The noise-driven iterates}}
\author{LangevinLean Project}
\date{}
\begin{document}\maketitle
Draft mathematical exposition; these notes may contain mistakes.


Let \(U:E_d\to\mathbb R\) be a potential with gradient \(\nabla U\), let \(h\in\mathbb R\), and let \(T_h^z(x)=x-h\nabla U(x)+\sqrt{2h}\,z\) be the ULA update. Given an initial vector \(x_0\in E_d\) and a sequence of noise vectors \(\xi_0,\xi_1,\ldots\in E_d\), the ULA iterates are


\[X_0=x_0,\qquad X_{k+1}=T_h^{\xi_k}(X_k)=X_k-h\nabla U(X_k)+\sqrt{2h}\,\xi_k\qquad(k\ge0).\]


The iterate \(X_n\) depends only on \(x_0\) and \(\xi_0,\dots,\xi_{n-1}\). When \(x_0\) and the \(\xi_k\) are random vectors, with the \(\xi_k\) independent standard Gaussian vectors independent of \(x_0\), this is the unadjusted Langevin algorithm. The recursion itself is deterministic and is defined for every input sequence; the coupling estimates of the later sections hold for every common noise sequence. This section records the defining identities and proves that every iterate is a measurable function of the randomness when the gradient is Lipschitz and the inputs are measurable, so that the law of each iterate is defined.


\section{\texorpdfstring{Definition of the iterates}{Definition of the iterates}}


\textbf{Definition.} For \(U:E_d\to\mathbb R\), \(h\in\mathbb R\), \(x_0\in E_d\) and \((\xi_k)_{k\ge0}\) in \(E_d\), the ULA iterates are defined recursively by \(X_0=x_0\) and \(X_{k+1}=T_h^{\xi_k}(X_k)\) for \(k\ge0\).

\begin{leanbox}{Lean: \texttt{ulaChain}}
\leansource{\texttt{ulaChain} \quad UnadjustedLangevin/\allowbreak{}Dynamics.lean, lines 20–23}
\begin{leancode}
\leanline{0}{def\ ulaChain\ \{d\ :\ ℕ\}\ (U\ :\ E\ d\ →\ ℝ)\ (h\ :\ ℝ)\ (initial\ :\ E\ d)}
\leanline{4}{(noise\ :\ ℕ\ →\ E\ d)\ :\ ℕ\ →\ E\ d}
\leanline{2}{|\ 0\ =>\ initial}
\leanline{2}{|\ n\ +\ 1\ =>\ ulaUpdate\ U\ h\ (ulaChain\ U\ h\ initial\ noise\ n)\ (noise\ n)}
\end{leancode}
\leanlabel{How the Lean proof follows the argument}
\texttt{ula\allowbreak{}Chain U h initial noise} is a function \texttt{\ensuremath{\mathbb N} \ensuremath{\to} E d} defined by structural recursion: the clause for \texttt{0} is \texttt{initial} and the clause for \texttt{n + 1} is \texttt{ula\allowbreak{}Update U h (ula\allowbreak{}Chain U h initial noise n) (noise n)}. The noise is a sequence \texttt{noise : \ensuremath{\mathbb N} \ensuremath{\to} E d}.
\end{leanbox}

\section{\texorpdfstring{The initial iterate}{The initial iterate}}


\textbf{Claim.} \(X_0=x_0\).




\textbf{Proof.} This is the first clause of the definition. \(\square\)

\begin{leanbox}{Lean: \texttt{ulaChain\_\allowbreak{}zero}}
\leansource{\texttt{ulaChain\_\allowbreak{}zero} \quad UnadjustedLangevin/\allowbreak{}Dynamics.lean, lines 25–26}
\begin{leancode}
\leanline{0}{@[simp]\ theorem\ ulaChain\_zero\ \{d\ :\ ℕ\}\ (U\ :\ E\ d\ →\ ℝ)\ (h\ :\ ℝ)}
\leanline{4}{(initial\ :\ E\ d)\ (noise\ :\ ℕ\ →\ E\ d)\ :\ ulaChain\ U\ h\ initial\ noise\ 0\ =\ initial}
\end{leancode}
\leanlabel{How the Lean proof follows the argument}
\texttt{ula\allowbreak{}Chain\_\allowbreak{}zero} is proved by \texttt{rfl}; the attribute \texttt{@[simp]} makes it available to the simplifier.
\end{leanbox}

\section{\texorpdfstring{The update identity}{The update identity}}


\textbf{Claim.} For every \(k\ge0\), \(X_{k+1}=T_h^{\xi_k}(X_k)\).




\textbf{Proof.} This is the second clause of the definition. \(\square\)

\begin{leanbox}{Lean: \texttt{ulaChain\_\allowbreak{}succ}}
\leansource{\texttt{ulaChain\_\allowbreak{}succ} \quad UnadjustedLangevin/\allowbreak{}Dynamics.lean, lines 28–31}
\begin{leancode}
\leanline{0}{@[simp]\ theorem\ ulaChain\_succ\ \{d\ :\ ℕ\}\ (U\ :\ E\ d\ →\ ℝ)\ (h\ :\ ℝ)}
\leanline{4}{(initial\ :\ E\ d)\ (noise\ :\ ℕ\ →\ E\ d)\ (n\ :\ ℕ)\ :}
\leanline{4}{ulaChain\ U\ h\ initial\ noise\ (n\ +\ 1)\ =}
\leanline{6}{ulaUpdate\ U\ h\ (ulaChain\ U\ h\ initial\ noise\ n)\ (noise\ n)}
\end{leancode}
\leanlabel{How the Lean proof follows the argument}
\texttt{ula\allowbreak{}Chain\_\allowbreak{}succ} is proved by \texttt{rfl} and is a \texttt{@[simp]} lemma.
\end{leanbox}

\section{\texorpdfstring{Measurability of every iterate}{Measurability of every iterate}}


\textbf{Claim.} Let \(\Omega\) be a measurable space, \(L\ge0\), and \(U:E_d\to\mathbb R\) with \(L\)-Lipschitz gradient. Let \(x_0:\Omega\to E_d\) and \(\xi_k:\Omega\to E_d\) (\(k\ge0\)) be measurable, where \(E_d\) carries its Borel \(\sigma\)-algebra. Then for every \(n\), the map \(\omega\mapsto X_n(\omega)\), the \(n\)-th iterate with initial vector \(x_0(\omega)\) and noise \((\xi_k(\omega))_k\), is measurable.




\textbf{Proof.} A Lipschitz map is continuous, and continuous maps between metric spaces are Borel measurable; so \(\nabla U\) is measurable. We argue by induction on \(n\). For \(n=0\), \(X_0=x_0\) is measurable by hypothesis. Assume \(X_n\) is measurable. Then \(\nabla U\circ X_n\) is measurable as a composition of measurable maps. Scalar multiples, differences and sums of measurable \(E_d\)-valued maps are measurable (the vector operations are continuous and \(E_d\) is separable, so the Borel \(\sigma\)-algebra of \(E_d\times E_d\) is the product \(\sigma\)-algebra). Hence \(X_{n+1}=X_n-h\,\nabla U\circ X_n+\sqrt{2h}\,\xi_n\) is measurable. \(\square\)




Consequently the law \((X_n)_\#P\) of every iterate is defined for every measure \(P\) on \(\Omega\).

\begin{leanbox}{Lean: \texttt{ulaChain\_\allowbreak{}measurable}}
\leansource{\texttt{ulaChain\_\allowbreak{}measurable} \quad UnadjustedLangevin/\allowbreak{}Dynamics.lean, lines 34–46}
\begin{leancode}
\leanline{0}{theorem\ ulaChain\_measurable\ \{Ω\ :\ Type*\}\ [MeasurableSpace\ Ω]\ \{d\ :\ ℕ\}}
\leanline{4}{(U\ :\ E\ d\ →\ ℝ)\ (h\ :\ ℝ)\ (L\ :\ ℝ≥0)}
\leanline{4}{(hLip\ :\ LipschitzWith\ L\ (gradient\ U))}
\leanline{4}{(initial\ :\ Ω\ →\ E\ d)\ (noise\ :\ ℕ\ →\ Ω\ →\ E\ d)}
\leanline{4}{(hinitial\ :\ Measurable\ initial)\ (hnoise\ :\ ∀\ n,\ Measurable\ (noise\ n))\ (n\ :\ ℕ)\ :}
\leanline{4}{Measurable\ (fun\ ω\ ↦\ ulaChain\ U\ h\ (initial\ ω)\ (fun\ k\ ↦\ noise\ k\ ω)\ n)}
\end{leancode}
\leanlabel{How the Lean proof follows the argument}
In \texttt{ula\allowbreak{}Chain\_\allowbreak{}measurable} the hypotheses are \texttt{h\allowbreak{}Lip : Lipschitz\allowbreak{}With L (gradient U)}, \texttt{hinitial : Measurable initial} and \texttt{hnoise : \ensuremath{\forall} n, Measurable (noise n)}. \texttt{hg} is measurability of the gradient, \texttt{h\allowbreak{}Lip.\allowbreak{}continuous.\allowbreak{}measurable}. The proof is \texttt{induction n}: the \texttt{zero} case is \texttt{hinitial}; in the \texttt{succ} case \texttt{simp only [ula\allowbreak{}Chain\_\allowbreak{}succ, ula\allowbreak{}Update]} unfolds one step, and the term \texttt{(ih.\allowbreak{}sub ((hg.\allowbreak{}comp ih).\allowbreak{}const\_\allowbreak{}smul h)).\allowbreak{}add ((hnoise n).\allowbreak{}const\_\allowbreak{}smul (Real.\allowbreak{}sqrt (2 * h)))} assembles measurability from the induction hypothesis \texttt{ih} by the rules for composition, scalar multiplication, subtraction and addition.
\end{leanbox}

\end{document}
